\documentclass[11pt]{article}

% --------------------------------------------------
% Packages
% --------------------------------------------------
\usepackage[T1]{fontenc}
\usepackage{lmodern}
\usepackage{geometry}
\usepackage{setspace}
\usepackage{amsmath, amssymb, amsthm}
\usepackage{csquotes}
\usepackage{hyperref}

\geometry{margin=1in}
\setstretch{1.15}

% --------------------------------------------------
% Theorem environments
% --------------------------------------------------
\newtheorem{definition}{Definition}
\newtheorem{proposition}{Proposition}
\newtheorem{theorem}{Theorem}

% --------------------------------------------------
% Title information
% --------------------------------------------------
\title{Neural Commit Semantics and the Algebraic Structure of Language}
\author{Flyxion}
\date{December 2025}

\begin{document}

\maketitle

\begin{abstract}
Human language exhibits algebraic properties—most notably commutativity with respect to linear order and nonassociativity under hierarchical composition—that resist explanation in purely statistical or state-based models. Recent neurocomputational theories, such as Murphy’s ROSE framework, suggest that these properties are enforced by oscillatory control dynamics rather than symbolic rule encoding. In this paper, we provide a formal account of this claim by introducing \emph{neural commit semantics}: an event-first architecture in which irreversible commit operations enforce algebraic invariants over replayable histories. We show that commutativity arises naturally from order-insensitive merge under canonical collapse, while nonassociativity is enforced by commit-induced irreversibility. We reinterpret ROSE as a perceptual control system operating over an authoritative event log, and we prove that contemporary large language models lack the architectural capacity to enforce these invariants under counterfactual intervention. The result is a unified account of linguistic structure, neural dynamics, and explanation itself, grounded in counterfactual control rather than prediction.
\end{abstract}

\section{Introduction}

Human language exhibits a combination of expressive power and structural constraint that remains difficult to explain within prevailing models of cognition and computation. Linguistic expressions are unbounded in principle, yet tightly constrained in form; they permit flexible recombination while enforcing strict algebraic invariants. Two such invariants are particularly striking: syntactic composition is commutative with respect to linear order, yet nonassociative with respect to hierarchical grouping. These properties are central to generative grammar, but they pose a challenge for models that learn from surface statistics alone.

Recent advances in large language models (LLMs) have intensified this challenge. Despite their success at predicting linguistic continuations and correlating with neural data, such models systematically fail to respect higher-order structural constraints, particularly under counterfactual manipulation. This raises a fundamental question: what kind of architecture is required not merely to predict linguistic behavior, but to \emph{explain} the structural invariants that define it?

A growing body of theoretical work suggests that the answer lies not in richer representations or larger datasets, but in a different notion of computation altogether. In particular, Murphy’s ROSE (Representation, Operation, Structure, Encoding) framework proposes that linguistic structure is enforced by neural control dynamics operating across multiple oscillatory timescales. On this view, algebraic properties such as headedness and nonassociativity are not stipulated symbolically, but emerge from the timing and coordination of neural activity—specifically through commit-like events mediated by oscillatory phase transitions.

In this paper, we provide a formal account of this intuition. We argue that linguistic algebra is enforced by \emph{neural commit semantics}: irreversible operations that collapse speculative compositions into indivisible units, thereby constraining subsequent computation. To make this precise, we adopt an event-first ontology in which semantic structure arises from replayable histories rather than mutable states. Within this framework, explanation is identified not with prediction or correlation, but with the capacity to evaluate counterfactual interventions—determining which structural manipulations are admissible and which are not.

Our contributions are threefold. First, we show that commutativity and nonassociativity follow necessarily from commit-based control over event composition, without requiring symbolic rules or learned constraints. Second, we reinterpret the ROSE architecture as a form of perceptual control operating over an authoritative event history, providing a principled link between linguistic theory and neural dynamics. Third, we prove that contemporary statistical sequence models lack the architectural resources to enforce these invariants, regardless of scale or training data.

The result is a unified account of linguistic structure, neural implementation, and explanation itself. Language, on this view, is not explained by prediction, representation, or optimization, but by the maintenance of algebraic invariants under counterfactual control.

\section{Commit Semantics and Linguistic Algebra}
\label{sec:commit_algebra}

This section presents the core technical claim of the paper: the algebraic properties of human syntax—commutativity with respect to linear order and nonassociativity with respect to hierarchical composition—are enforced by \emph{commit semantics} rather than by symbolic rules or learned statistical regularities. We show that once linguistic composition is understood as event-based and commit-governed, these properties follow necessarily from architectural constraints.

\subsection{The Algebraic Problem}

Two algebraic features of syntactic composition have long been emphasized in generative linguistics.

First, syntactic combination is insensitive to linear order. While word order may affect interpretation at the interfaces, the basic combinatorial operation does not depend on whether an element is encountered first or second. Second, syntactic composition is nonassociative: the grouping of constituents matters. The structure resulting from combining $x$ and $y$ before $z$ is not equivalent to that obtained by combining $y$ and $z$ before $x$.

Formally, these properties can be stated as follows. Let $\mathcal{M}$ denote the basic combinatorial operation (traditionally, Merge). Then:
\begin{enumerate}
  \item \emph{Commutativity}: $\mathcal{M}(x,y) \cong \mathcal{M}(y,x)$,
  \item \emph{Nonassociativity}: $\mathcal{M}(\mathcal{M}(x,y),z) \not\cong \mathcal{M}(x,\mathcal{M}(y,z))$.
\end{enumerate}

Standard symbolic approaches stipulate these properties axiomatically. Statistical sequence models, by contrast, typically violate both: they are order-sensitive by construction and freely reassociate internal representations. What is required is an architecture in which commutativity and nonassociativity are \emph{enforced}, not declared.

\subsection{Event-Based Composition}

We adopt an event-first ontology in which linguistic computation unfolds as a sequence of compositional events rather than transitions between stored states. Basic linguistic elements are treated as event-regions rather than symbols, and composition corresponds to the formation of larger regions.

\begin{definition}[Event-Based Merge]
Given two event-regions $x$ and $y$, merge is defined as region union:
\[
\mathsf{Merge}(x,y) \;=\; x \cup y.
\]
\end{definition}

Crucially, region union is symmetric. At this stage of computation, there is no privileged order of combination and no commitment to internal structure. Merge is therefore commutative by construction.

\subsection{Commit as Irreversible Collapse}

Unrestricted region union alone is insufficient to produce stable structure. Without further constraint, regions remain indefinitely revisable, and reassociation is always possible. The key architectural addition is \emph{commit}: an irreversible operation that collapses a composite region into an indivisible unit.

\begin{definition}[Commit]
A commit operation applies a canonical collapse operator $D$ to a composite region $R$, yielding a new atomic unit $h$:
\[
h \;=\; D(R).
\]
After commit, the internal structure of $R$ is no longer accessible to subsequent compositional operations.
\end{definition}

Commit is not a representational update but a control event. It marks the transition from speculative composition to stabilized structure.

\subsection{Enforcing Nonassociativity}

The irreversibility of commit has immediate algebraic consequences.

\begin{proposition}[Commit-Enforced Nonassociativity]
Once a composite region has been committed, reassociation of its internal constituents is structurally impossible.
\end{proposition}

\begin{proof}
Consider the two compositions:
\[
D(D(x \cup y) \cup z) \quad\text{and}\quad D(x \cup D(y \cup z)).
\]
In the first case, $x$ and $y$ are merged and committed before $z$ is introduced, yielding an atomic unit $h = D(x \cup y)$. Subsequent composition treats $h$ as indivisible. In the second case, $y$ and $z$ are committed first. Because commit renders internal structure inaccessible, there exists no sequence of operations that transforms one result into the other without violating irreversibility. Hence the two structures are not equivalent.
\end{proof}

Nonassociativity is therefore not imposed by grammatical fiat, but enforced by the timing of irreversible commit events.

\subsection{Headedness Without Order}

A further consequence of commit semantics concerns headedness. In symbolic theories, heads are often determined by positional or feature-based rules. In a commit-based architecture, headedness arises from dominance relations evaluated at commit time.

We assume that each event-region carries scalar dominance measures reflecting activation strength, stability, or coherence. At commit, the canonical representative is selected based on these measures.

Because dominance is invariant under permutation of inputs, headedness is independent of linear order. This yields commutativity “for free”: the same composite is obtained regardless of presentation order, provided the same dominance relations hold.

\subsection{Neural Interpretation}

Commit semantics admits a natural neural interpretation. Speculative merge operations correspond to transient local activity, while commit corresponds to a global coordination event that stabilizes structure. Oscillatory dynamics provide a plausible implementation: high-frequency activity supports proposal generation, slower rhythms define a workspace, and brief coordination bursts signal commitment and transition to the next compositional cycle.

On this view, algebraic structure is enforced by control flow rather than encoded representation. Neural dynamics do not store trees; they regulate when structure becomes fixed.

\subsection{Summary}

Commit semantics resolves a central explanatory problem in the theory of language. By distinguishing speculative composition from irreversible stabilization, it explains how syntactic composition can be simultaneously commutative and nonassociative. These properties emerge from architectural constraints on event histories, not from symbolic stipulation or statistical approximation. In the following section, we show how this framework integrates naturally with Murphy’s ROSE architecture and provides a principled basis for its neural implementation claims.

\section{ROSE as Neural Perceptual Control}
\label{sec:rose_control}

Having established how commit semantics enforces the algebraic structure of language, we now turn to the question of neural implementation. In this section, we reinterpret Murphy’s ROSE (Representation, Operation, Structure, Encoding) architecture not as a pipeline of representational stages, but as a form of \emph{perceptual control} operating over event-based composition. This reinterpretation preserves the empirical and theoretical motivations of ROSE while situating it within a counterfactually complete, event-first framework.

\subsection{From Processing Stages to Control Roles}

ROSE is often presented as a decomposition of linguistic computation into four components: representations (R), operations (O), structure (S), and encoding (E). Read naively, this decomposition invites a stage-based interpretation, in which information flows sequentially from representations to operations, structures, and finally encodings. Such readings risk reintroducing state-based assumptions that the ROSE framework was explicitly designed to avoid.

Instead, we propose that R, O, S, and E should be understood as \emph{control roles} rather than processing stages. Each role corresponds to a distinct way of probing, proposing, or stabilizing structure over a shared event history. All four operate concurrently over the same underlying compositional process, differing only in temporal scale and functional scope.

\subsection{R and O: Probing and Proposing}

The R-component corresponds to the detection of atomic features relevant for linguistic composition. In an event-based framework, these are not stored representations but transient indices into the current history of events. They provide access to what is available for composition, not a persistent encoding of content.

The O-component corresponds to speculative compositional activity: candidate merge operations generated in response to available features. Importantly, these operations are \emph{proposals}, not commitments. They explore possible ways of combining event-regions without yet stabilizing structure. This aligns naturally with Murphy’s characterization of O as operating at fast timescales and being closely associated with local, high-frequency neural activity.

Crucially, neither R nor O has the authority to fix structure. Their role is exploratory and provisional, supplying candidates to a higher-level control mechanism.

\subsection{S: Structural Invariants as Controlled Variables}

The S-component is often described as encoding hierarchical structure. Within a control-theoretic interpretation, its role is more precise: S specifies the \emph{invariants} that must be maintained across compositional cycles. These invariants include hierarchical depth, headedness relations, and the integrity of previously committed units.

Rather than constructing trees directly, S monitors whether speculative operations are compatible with existing commitments. In this sense, S functions as the reference signal in a perceptual control loop: it defines what counts as an acceptable structural outcome, against which candidate operations are evaluated.

\subsection{E: Encoding as Coordination, Not Representation}

The E-component mediates large-scale coordination across neural systems. In traditional accounts, this might be taken to involve representational encoding suitable for memory or interface with other cognitive systems. In the present framework, E instead serves to synchronize control across distributed regions.

Encoding, on this view, is not about storing structure but about ensuring that committed structure is consistently respected across different neural populations and timescales. This interpretation aligns with Murphy’s emphasis on inter-areal dynamics and traveling waves, which coordinate activity without localizing structure to a single site.

\subsection{Commit as the Control Pivot}

Within this reinterpretation, commit semantics occupies a central position. Commit events mark the moment at which speculative operations are evaluated against structural invariants and either stabilized or rejected. They serve as the pivot of the control loop, transforming provisional composition into authoritative structure.

Neurally, this corresponds to brief, system-level coordination events—such as beta-mediated bursts—that silence ongoing speculation, consolidate the result of the current compositional cycle, and prepare the system for the next. The critical point is that these events enforce algebraic constraints by \emph{terminating degrees of freedom}. Once committed, internal structure is no longer revisable, ensuring nonassociativity and stability.

\subsection{ROSE as Perceptual Control}

Taken together, these roles form a closed-loop control system:
\begin{itemize}
  \item R provides access to available elements,
  \item O generates candidate compositions,
  \item S specifies the invariants to be maintained,
  \item E coordinates stabilization across the system,
  \item commit events enforce irreversible decisions.
\end{itemize}

This architecture is naturally described as perceptual control. The system does not build representations of syntax in order to manipulate them; instead, it acts to maintain structural invariants over time in the face of noisy input and competing proposals. Errors correspond not to mismatches between representations, but to violations of admissibility that trigger repair or reanalysis.

\subsection{Implications}

Recasting ROSE as a control architecture clarifies several points of contention. It explains how abstract algebraic properties can be neurally enforced without symbolic encoding, why linguistic computation is distributed rather than localized, and why oscillatory dynamics play a central role. It also clarifies the limits of purely statistical models: without commit semantics and invariant-based control, they lack the means to enforce structure rather than merely approximate it.

In the next section, we leverage this control-theoretic perspective to show why contemporary large language models, despite their empirical success, are architecturally incapable of enforcing the algebraic constraints that define human language.

\section{Why Large Language Models Fail}
\label{sec:llm_failure}

The preceding sections identified commit semantics and invariant-based control as the architectural basis of linguistic structure. We now show that contemporary large language models (LLMs), despite their empirical success, are incapable of enforcing these constraints. The limitation is not a matter of scale, training data, or objective function, but of fundamental architectural mismatch.

\subsection{Prediction Without Authority}

LLMs are trained to optimize the conditional probability of token sequences. Their internal computations operate over continuously updated activation states, and all structure remains revisable until output is produced. At no point does the model enforce an authoritative distinction between speculative composition and committed structure.

This absence of authority has immediate consequences. Because all internal representations remain accessible, any part of a previously composed structure may be reweighted or reinterpreted in light of subsequent input. The model therefore lacks a mechanism for irreversible commitment.

\subsection{No Commit, No Nonassociativity}

As shown in Section~\ref{sec:commit_algebra}, nonassociativity requires that once a composite is formed, its internal constituents become inaccessible for reassociation. LLMs violate this requirement by design.

\begin{proposition}[Absence of Commit Semantics in LLMs]
LLMs cannot enforce nonassociativity under counterfactual intervention.
\end{proposition}

\begin{proof}
Transformer architectures maintain simultaneous access to all token representations at every layer via attention. No operation collapses a subset of representations into an indivisible unit whose internal structure is no longer addressable. Consequently, for any composition of elements $x,y,z$, the model can freely approximate both $(x \cdot y) \cdot z$ and $x \cdot (y \cdot z)$ by reweighting attention, even if one structure is linguistically illicit. Thus nonassociativity is not enforced.
\end{proof}

\subsection{Order Sensitivity and Spurious Commutativity}

Although LLMs may appear insensitive to linear order in some contexts, this behavior is contingent and statistical. Because token order is explicitly encoded in positional embeddings, commutativity is not a structural invariant of the model.

Any apparent order-independence reflects properties of the training distribution rather than architectural enforcement. Under counterfactual perturbation—such as adversarial reordering—these models readily violate commutativity constraints that human speakers respect.

\subsection{Structural Ambiguity Versus Uncertainty}

A central distinction in linguistic theory is between \emph{structural ambiguity} and mere uncertainty. Structural ambiguity arises when a single surface string corresponds to multiple distinct underlying structures, each internally coherent. Uncertainty arises when a system lacks sufficient information to decide among alternatives.

LLMs collapse this distinction. Because they operate only over surface sequences, alternative parses are represented as probability mass over continuations, not as distinct admissible histories. As a result, the model cannot represent the fact that multiple structures are equally valid while others are categorically excluded.

\subsection{Correlation Is Not Explanation}

It is sometimes argued that high correlation between LLM activations and neural data suffices for explanatory relevance. From the perspective developed here, this is a category error. Correlation reflects similarity in surface projections; explanation requires counterfactual competence.

A system explains linguistic structure only if it can determine whether a proposed intervention—such as reassociating constituents or violating headedness—is admissible. LLMs lack this capacity. They can predict likely continuations, but they cannot reject structurally impossible ones except insofar as such cases are rare in the training data.

\subsection{Why Scaling Does Not Help}

The limitations described above are invariant under scale. Increasing model size or data improves approximation of surface distributions, but it does not introduce irreversible commit operations or authoritative event histories. Without these, the algebraic invariants of language remain unenforced.

This explains why failures of higher-order generalization persist even in the largest models: the relevant constraints are architectural, not statistical.

\subsection{Summary}

Large language models succeed by approximating the projections of linguistic structure. They fail because they lack the control architecture required to enforce structure itself. Without commit semantics, invariant-based control, and counterfactual admissibility, no amount of training can transform prediction into explanation. The contrast with human language is therefore not one of degree, but of kind.


\begin{thebibliography}{99}

\bibitem{murphy2023rose}
Murphy, E. (2023).
\newblock \emph{ROSE: A Neurocomputational Architecture for Syntax}.
\newblock arXiv preprint arXiv:2303.08877.
\newblock Available at \texttt{https://arxiv.org/abs/2303.08877}.

\bibitem{chomsky1995minimalist}
Chomsky, N. (1995).
\newblock \emph{The Minimalist Program}.
\newblock MIT Press.

\bibitem{chomsky2016what}
Chomsky, N. (2016).
\newblock What kind of creatures are we?
\newblock \emph{Columbia University Press}.

\bibitem{pearl2009causality}
Pearl, J. (2009).
\newblock \emph{Causality: Models, Reasoning, and Inference}.
\newblock Cambridge University Press.

\bibitem{pearl2018book}
Pearl, J., \& Mackenzie, D. (2018).
\newblock \emph{The Book of Why}.
\newblock Basic Books.

\bibitem{friston2010free}
Friston, K. (2010).
\newblock The free-energy principle: a unified brain theory?
\newblock \emph{Nature Reviews Neuroscience}, 11(2), 127--138.

\bibitem{levin2022tame}
Levin, M. (2022).
\newblock Technological approach to mind everywhere.
\newblock \emph{Trends in Cognitive Sciences}, 26(9), 766--784.

\bibitem{marr1982vision}
Marr, D. (1982).
\newblock \emph{Vision: A Computational Investigation into the Human Representation and Processing of Visual Information}.
\newblock W. H. Freeman.

\bibitem{vaswani2017attention}
Vaswani, A., Shazeer, N., Parmar, N., et al. (2017).
\newblock Attention is all you need.
\newblock \emph{Advances in Neural Information Processing Systems}, 30.

\bibitem{lake2017building}
Lake, B. M., Ullman, T. D., Tenenbaum, J. B., \& Gershman, S. J. (2017).
\newblock Building machines that learn and think like people.
\newblock \emph{Behavioral and Brain Sciences}, 40, e253.

\bibitem{newell1976computer}
Newell, A. (1976).
\newblock Computer science as empirical inquiry.
\newblock \emph{Communications of the ACM}, 19(3), 113--126.

\end{thebibliography}

\end{document}
